% Part: counterfactuals
% Chapter: minimal-change-semantics
% Section: transitivity

\documentclass[../../../include/open-logic-section]{subfiles}

\begin{document}

\olfileid[hi]{cnt}{min}{cpo}

\olsection{प्रतिधनात्मकता}

भौतिक और कठोर सशर्त अपने प्रतिधनात्मक के समतुल्य होते हैं। प्रतितथ्यात्मक
सशर्त ऐसे नहीं होते। क्राट्ज़र का यह उदाहरण देखें:
\begin{quote}
  यदि ग्योथे 1832 में न मरे होते, तो वे अब (फिर भी) मृत होते।

  यदि ग्योथे अब मृत न होते, तो वे 1832 में मरे होते।
\end{quote}
पहला वाक्य सत्य है: मनुष्य सैकड़ों वर्ष नहीं जीते। दूसरा स्पष्टतः असत्य
है: यदि ग्योथे अब मृत न होते, तो वे अभी जीवित होते और इसलिए 1832 में मर
नहीं सकते थे।

\begin{ex}\ollabel{ex:contraposition-counterex}
  गोलक अर्थविज्ञान प्रतिधनात्मकता को अमान्य करता है; अर्थात हमारे पास $p
  \cif q \Entails/ \lnot q \cif \lnot p$ है। $p$ को ``ग्योथे 1832 में
  नहीं मरे'' और $q$ को ``ग्योथे अब मृत हैं'' मानें। इसे हम ऐसे प्रतिरूप
  $\mModel{M_1} = \tuple{W, O, V}$ में निरूपित कर सकते हैं जिसमें $W =
  \{w, w_1, w_2\}$, $O = \{\{w\}, \{w, w_1\}, \{w, w_1, w_2\}\}$,
  $V(p) = \{w_1, w_2\}$ और $V(q) = \{w, w_1\}$। अतः $w$ वास्तविक
  जगत है, जहाँ ग्योथे 1832 में मरे और अब भी मृत हैं; $w_1$ वह (निकट)
  जगत है, जहाँ ग्योथे, मान लें, 1833 में मरे और अब भी मृत हैं; और $w_2$
  वह (दूरस्थ) जगत है, जहाँ ग्योथे अब भी जीवित हैं।
\begin{figure}
\begin{center}
\begin{tikzpicture}[scale=.6,layerwidth=1.5]\tiny
  \spheresystem{3}
  \begin{scope}
    \clip \propositionplot{0}{.8}{50}{4.5} -- (4.5,-4.5) -- (-4.5,-4.5) -- (-4.5,4.5) -- (4.5,4.5);
    \spherelayer{3}
  \end{scope}
  \propositionintersect{0}{2}{110}{5.5}
  \proposition{0}{.8}{50}{4.5}
  \path (0,0) node[world] {$w$};
  \spherepos{0}{2}{node[world] {$w_1$}}
  \spherepos{-50}{3}{node[world] {$w_2$}}
  \spherepos{28}{3}{node {$q$}}
  \spherepos{42}{3}{node {$\lnot q$}}
  \spherepos{-55}{4}{node {$p$}}
  \spherepos{-67}{4}{node {$\lnot p$}}
\end{tikzpicture}
\caption{प्रतिधनात्मकता का प्रतिदृष्टांत}
\ollabel{fig:contraposition}
\end{center}
\end{figure}
एक $p$-स्वीकारी गोलक $S = \{w, w_1\}$ है और उसके सभी जगतों में
$p \lif q$ सत्य है, अतः $\mSat{M}{p \cif q}[w]$। किंतु $\lnot q$-स्वीकारी
गोलक $\{w, w_1, w_2\}$ में एक जगत, अर्थात~$w_2$, ऐसा है जहाँ $q$ असत्य
और $p$ सत्य है, अतः $\mSat/{M}{\lnot q \lif \lnot p}[w_2]$।
\end{ex}

\end{document}
